\section{Related Work}%
\label{sec:related}

\paragraph{Bunched typing and bunched implication.}

The idea of expressing fine-grained resource control within a program
logic as introduced in Bunched Implication~\cite{OHP99} is central to
modern-day program verifiers~\cite{AD+14,Iris,SLK+21,MSJ16,GM+25}
based on separation logic~\cite{SeparationLogic,OH19}.  Transplanting
the notion of \emph{bunched typing} as introduced
by~\citet{BunchedType}, using $\ctcomma$ and $\ctstar$ to track
entangled and independent resource use within a type context has
received much less attention.  Recently proposed type systems and
type-based verification frameworks that support resource
reasoning~\cite{FR+21,PLM+23,JJKD18} do so by exposing an explicit
ownership discipline, rather than manipulating a bunched context
structure.  Notably, these other systems are designed primarily with
safety in mind, and do not consider the role of resources and contexts
as a unifying mechanism for safety and reachability.

\paragraph{Reachability verification.}
Recent work on reachability verification in type systems includes
two-sided type systems~\cite{Ill-types} which refute typing formulas
to show that ill-formed programs do not evaluate and a complementary
approach~\cite{RamsayIncorType} that internalizes negation directly as
a complement operator on types to support incorrectness reasoning.
Coverage types~\cite{CoverageType} underapproximate
test-input-generator behavior to establish coverage completeness.
While these efforts emphasize an algorithmic characterization of
reachability within a type system, they do not consider the logical
underpinning that relates reachability and safety claims as we do
here, and instead treat reachability properties independently of other
safety considerations.

We note that contemporary work on \emph{Reachability
  Types}~\cite{ReachabilityTypes, WBJBR+24} uses ``reachability'' in a
sense unrelated to how it is used here, and focuses on soundly
tracking aliasing of mutable references in higher-order effectful
programs.

\paragraph{Incorrectness Logics.}
Our angelic modality can be viewed as establishing a witness to a
property, and thus is closely related to operators found in
underapproximate logics~\cite{ReverseHoareLogic, OHP19, RBD+20,
  RHLE, HyperHoareLogic, QuantitativeWeakestHyperPre, UTurn,
  BackwardIncorrectness} that establish judgements which guarantee the
existence of certain executions rather than ensuring the absence of
unsafe executions.

However, these logics and their descendents either fix a single
approximation direction per triple~\cite{ReverseHoareLogic, OHP19,
  RBD+20} or disallow approximation entirely~\cite{ESL+23}, and do not
support the construction or decomposition of context structure, which
is a central feature of our type system. A complementary line of work
has studied the relationship between various correctness and
incorrectness program logics, when viewed through the lens of
predicate transformers and Galois connections~\cite{Cousot+24, VK+25};
these efforts relate distinct program logics to one another
externally, whereas our type system internalizes both approximation
directions into a single reasoning framework.

Like us, outcome logic~\cite{ZDS23, Zil25TOPLAS} also unifies safety
and reachability through branch composition $\varphi \oplus \psi$, and
its $\oplus$ adjoins control-flow alternatives like our $\ctoplus$.
But, mixing angelic and demonic constraints within a single outcome
logic judgement requires an \emph{a priori} partitioning of outcomes
via $\oplus$, whereas our type system determines valid angelic choices
compositionally, without requiring such a partition to be specified
upfront. Outcome logic has since been extended with, e.g.,
separation-logic-style local reasoning~\cite{ZSS+24} and to programs
featuring adversarially-resolved nondeterminism~\cite{ZKST+25}. In
these extensions, however, 'demonic' and 'angelic' characterize how a
program's nondeterminism is resolved by its semantics, whereas our
modalities are specification-level annotations that determine whether
an individual refinement is read as an upper or lower bound, and may
be mixed freely within a single type judgment.

\paragraph{May- and must-style program analysis.}
There has been significant prior work on combining may- and must-style
reasoning within a program analysis~\cite{GHK+06,GN+10,NYC+12,BVR+08}.
For example,~\citet{GN+10} alternates overapproximate may-summaries
with underapproximate must-summaries so that callee information can
inform caller analysis.

Refutation is a type-level specialization of may/must reasoning: a
failed refinement type check (which is may-style) does not by itself
establish that a bug is reachable (which requires must-style
evidence).  Refinement type
refutations~\cite{RefinementTypeRefutations} address this gap by
deriving \emph{must}-style witnesses that explain why typing failed,
often through concrete counterexample instantiations, while
symbolic-execution tooling~\cite{RefinementTypeSymbolicExecution}
searches for such counterexamples externally.  More broadly, our
context types provide a unified theoretical foundation for both
approaches by exploiting the duality between safety and reachability
reasoning.

\paragraph{Context-aware typing}
Coeffect and graded modal type systems~\cite{Coeffects, Katsumata+14,
  GaboardiCombining, OrchardQuantitative} annotate programs with
context information such as usage, liveness, and security levels,
interpreted via graded (co)monads.  Like our context types, they treat
contextual information as a typed resource with a compositional
algebra, but they track the demands a computation makes on its
environment or what effects the computation has on its environment,
rather than how a typing context \emph{interprets} value sets under
demonic or angelic modalities.  Thus, they do not provide a mechanism
to internalize safety/reachability modalities on the same base
refinement, nor do they provide a bunched capability structure for
tracking dependencies between bindings in a higher-order setting.
